\documentclass{article}

% TMLR format
\usepackage[accepted]{tmlr}

\newcommand{\openreview}{\url{https://openreview.net/forum?id=XXXX}}
\def\month{06}
\def\year{2026}

\usepackage{amsmath,amssymb,amsthm}
\usepackage{booktabs}
\usepackage{graphicx}
\usepackage{hyperref}
\usepackage{xcolor}
\usepackage{array}
\usepackage{enumitem}
\usepackage{multirow}

\newcommand{\Gr}{\mathrm{Gr}}
\newcommand{\IIA}{\mathrm{IIA}}
\newcommand{\R}{\mathbb{R}}
\newcommand{\Z}{\mathbb{Z}}
\newcommand{\dgr}{d_{\Gr}}
\newcommand{\Span}{\mathrm{span}}
\newcommand{\KL}{D_\mathrm{KL}}
\newcommand{\JS}{D_\mathrm{JS}}
\DeclareMathOperator{\diag}{diag}

\newtheorem{assumption}{Assumption}
\newtheorem{proposition}{Proposition}

\title{When Does Linear Causal Abstraction Work?\\
Mapping the Boundary on the Grassmannian}

\author{
Elliot Tower \\
\texttt{elliot@elliottower.com}
}

\begin{document}
\maketitle

\begin{abstract}
Distributed Alignment Search (DAS) identifies low-dimensional linear subspaces
that mediate model behavior under interchange interventions, implicitly assuming
that the relevant causal variable lives on the Grassmannian manifold
$\Gr(k,d)$. When does this assumption hold, and what happens when it fails?

We construct an atlas of 14 modular arithmetic operations spanning the boundary
between linear and nonlinear causal structure. Three findings emerge.
\textbf{(1)}~Operations partition into three sharp classes: those that
\emph{always} produce Grassmannian causal variables across random seeds, those
where the outcome is \emph{stochastic}---same operation, different seeds,
opposite causal geometries---and those that \emph{never} do. The boundary is
governed by grokking: Grassmannian variables appear if and only if the model
generalizes. \textbf{(2)}~Linear DAS returns \emph{exactly zero} IIA on fully
grokked modular addition at $k \leq 16$, confirming that the causal variable is
fundamentally nonlinear despite the model having learned the correct algorithm.
Memorized models achieve $\IIA \geq 0.86$ at $k = 2$ without any geometric
structure, demonstrating that IIA alone cannot distinguish genuine causal
variables from lookup tables. \textbf{(3)}~We introduce a Structured pi-SAE---a label-conditional sparse
autoencoder that disentangles causal from nuisance factors---recovering
structured causal variables for operations where linear DAS fails entirely.
With end-to-end intervention training, the Structured pi-SAE achieves
IIA~=~0.97 on GPT-2 hypernymy at $k{=}1$, where DAS gets 0.07.
For grokked operations, the latent space shows equivariance exceeding 95\%.
For never-Grassmannian operations, it reveals the \emph{nonlinear} causal
structure that DAS cannot access.

These results establish a hierarchy: IIA is necessary but insufficient;
equivariance distinguishes genuine structure from memorization; and structured
nonlinear methods extend causal variable discovery beyond the Grassmannian.
\textbf{(4)}~We show that unconstrained nonlinear methods (MLP featurizer +
linear DAS) achieve perfect IIA by learning degenerate encoder-decoders---lookup
tables that ignore the base input entirely. We introduce \emph{intervention
faithfulness} metrics---diversity ratio, reconstruction fidelity, and continuous
distributional measures (KL, Jensen-Shannon)---that expose this failure mode.
The structured VAE's ELBO prevents degeneracy, establishing that generative
structure is necessary for valid nonlinear causal abstraction.
\end{abstract}

\section{Introduction}
\label{sec:intro}

Mechanistic interpretability seeks to identify causally active,
human-understandable variables inside neural networks. Causal abstraction
formalizes this goal: a model variable is ``causal'' if interventions on its
internal representation behave like interventions on a high-level algorithmic
variable \citep{geiger2021causal}. Distributed Alignment Search (DAS)
operationalizes causal abstraction by learning a rotation $Q \in \R^{d \times k}$
that identifies a $k$-dimensional subspace of the residual stream. Swapping the
in-subspace component between two inputs and measuring whether the model
produces the corresponding counterfactual output yields Interchange Intervention
Accuracy (IIA) \citep{geiger2023finding}.

DAS has been applied to syntax \citep{geiger2023finding}, arithmetic
\citep{wu2024interpretability}, and board games \citep{li2023emergent}, but
every application embeds a structural assumption: the causal variable is a
\emph{linear} subspace---a point on the Grassmannian manifold $\Gr(k,d)$. When
the true causal variable is linear, DAS recovers it. When it is not, DAS may
still return a high-IIA subspace that reflects memorization or distributional
artifacts rather than genuine causal structure.

We do not ask whether DAS ``works'' in the narrow sense of achieving high IIA---it
almost always does. We ask: \textbf{when does the linear assumption hold?} Is the
subspace DAS recovers a genuine Grassmannian point with the geometric properties
expected of a linear causal variable, or an artifact of searching in the wrong
function class? And when it fails, \textbf{can a nonlinear method recover the
causal variable?}

To answer these questions, we construct an atlas of 14 operations over $\Z_p$
spanning the boundary between linear and nonlinear causal structure. For each, we
train a single-layer transformer in the standard grokking setup, fit DAS at
multiple subspace dimensions, and evaluate the recovered subspace with three
diagnostics beyond IIA: equivariance under the operation's symmetry group, circle
geometry of the DAS projection, and intrinsic dimension from $k$-sweeps. We then
apply a structured VAE \citep{narayanaswamy2017learning} to the
never-Grassmannian operations, showing that nonlinear causal variables exist even
where linear methods fail.

\paragraph{The linearity boundary is sharp and predictable.}
Operations partition into three classes. Seven \emph{always} produce Grassmannian
variables: equivariance exceeds 95\% and $k=2$ suffices for $\IIA \geq 0.96$.
Two are \emph{stochastic}: the same operation at different random seeds sometimes
produces Grassmannian structure and sometimes does not. The remaining five
\emph{never} produce Grassmannian variables. The mechanism is grokking: across all
operations and seeds, the model produces a Grassmannian variable if and only if it
generalizes.

\paragraph{Linear DAS fails completely on nonlinear causal variables.}
On fully grokked modular addition, linear DAS returns $\IIA = 0.0$ at every
checkpoint from epoch 500 through 25{,}000, at $k = 16$ (one-eighth of the
residual stream). The model has learned a correct algorithm---its test loss is
$< 10^{-6}$---but the causal variable encoding the modular sum lives on a
circle in activation space, not on a linear subspace. This is not a failure of
DAS optimization; it is a category error: searching $\Gr(k,d)$ when the
variable lives on $S^1 \hookrightarrow \R^d$.

\paragraph{Structured VAE recovers nonlinear causal structure.}
We adapt the semi-supervised deep generative model of
\citet{narayanaswamy2017learning}, which disentangles latent factors into a
supervised $z_\text{choice}$ (trained to predict the operation's output) and an
unsupervised $z_\text{other}$ (nuisance variance). The VAE's nonlinear encoder
maps activations to a structured latent space where the causal variable is
explicitly separated. For grokked operations, the VAE latent space exhibits the
same equivariance as DAS. For never-Grassmannian operations, the VAE reveals
structured nonlinear causal variables invisible to DAS.

% TODO: Add CPCA-init DAS hard-example experiments for grokking operations
% TODO: Report multi-seed stability (10 seeds per stochastic operation)

\paragraph{Contributions.}
\begin{enumerate}[leftmargin=*]
\item We construct an \textbf{atlas of 14 operations} that systematically maps
  the boundary of linear causal abstraction, organized by whether DAS recovers
  Grassmannian variables.
\item We discover \textbf{stochastic grokking}: for certain operations,
  whether grokking occurs is seed-dependent, and Grassmannian structure appears
  if and only if grokking happens. This provides the cleanest evidence that
  generalization---not architecture or loss landscape---is the mechanism enabling
  linear causal variables.
\item We show that \textbf{linear DAS returns zero IIA on grokked modular
  addition}, confirming that the causal variable is fundamentally nonlinear even
  when the model has learned a correct generalizing algorithm.
\item We develop \textbf{geometric diagnostics} beyond IIA---equivariance testing
  and $k$-sweep saturation profiles---that distinguish genuine Grassmannian
  variables from memorization artifacts.
\item We show that \textbf{memorization produces high IIA}: squaring and cubing
  achieve $\IIA \geq 0.86$ at $k=2$ and $\IIA = 1.0$ by $k=4$ without grokking.
\item We introduce a \textbf{Structured pi-SAE for nonlinear causal variable
  discovery}---a label-conditional sparse autoencoder with end-to-end
  intervention training---that recovers equivariant representations for
  operations where DAS fails entirely, and achieves IIA~=~0.97 on GPT-2
  language tasks at $k{=}1$.
\item We show that \textbf{unconstrained nonlinear featurizers achieve perfect
  IIA vacuously}: an MLP encoder-decoder trained end-to-end on interchange loss
  learns a lookup table (intervention diversity ratio $\approx 0$), not a
  genuine nonlinear coordinate system. The structured VAE's ELBO prevents this
  degeneracy.
\item We introduce \textbf{intervention faithfulness metrics}---diversity ratio,
  reconstruction fidelity, and continuous distributional measures---as necessary
  complements to IIA. We advocate that causal abstraction claims require both
  interchange effectiveness (IIA) and interchange faithfulness (the intervention
  modifies only the targeted variable).
\end{enumerate}


\section{Background}
\label{sec:background}

\subsection{Causal abstraction and DAS}

Let $h(x) \in \R^d$ be an intermediate activation for input $x$. DAS searches
for an orthonormal $Q \in \R^{d \times k}$ whose column span defines a causal
subspace. Given a base input $x_b$ and source input $x_s$, the interchange
intervention replaces the in-subspace component:
\begin{equation}
h' = h_b - QQ^\top h_b + QQ^\top h_s .
\label{eq:intervention}
\end{equation}
IIA measures the fraction of interventions for which the model's output matches
the target counterfactual. High IIA indicates causal sufficiency: the subspace
contains enough information to drive behavior under interchange. But sufficiency
is not validity. A subspace may score high IIA because it encodes a memorized
lookup table, because correlated directions redundantly encode the same label, or
because the intervention distribution is too easy.
\citet{makelov2024illusion} demonstrated this concretely: subspace activation
patching can produce interpretability illusions where dormant pathways inflate IIA
without reflecting genuine causal structure. Our central claim is that IIA must be
paired with geometric diagnostics that test whether the subspace is a genuine
linear causal variable rather than an artifact of the linear function class.

\subsection{The Grassmannian manifold}
\label{sec:grassmannian_background}

The Grassmannian $\Gr(k,d)$ is the set of all $k$-dimensional linear subspaces
of $\R^d$. A DAS solution $Q \in \R^{d \times k}$ represents a point
$[Q] = \Span(Q)$ on this manifold; any $QR$ for $R \in O(k)$ represents the
same point. The geodesic distance between two subspaces is determined by their
\emph{principal angles}: if $Q_1, Q_2$ have orthonormal columns, the singular
values of $Q_1^\top Q_2$ are $\cos \theta_i$, and the geodesic distance is
\begin{equation}
\dgr([Q_1],[Q_2]) = \left(\sum_i \theta_i^2\right)^{1/2}
\label{eq:geodesic}
\end{equation}
\citep{edelman1998geometry, absil2008optimization}. Every application of DAS is
implicitly a search on this manifold. The question ``does DAS work?'' is
really ``does the true causal variable live on $\Gr(k,d)$?'' When it does, DAS
recovers a meaningful point. When it does not, DAS still returns a point on
$\Gr(k,d)$, but that point may be an artifact.

\subsection{Structured disentangled generative models}
\label{sec:structured_vae}

Semi-supervised deep generative models \citep{narayanaswamy2017learning} extend
the VAE framework \citep{kingma2014auto} by partitioning the latent space into
supervised and unsupervised factors. The generative model factors as
$p(x, y, z) = p(x \mid y, z) \, p(y) \, p(z)$, where $y$ is a labeled
(interpretable) variable and $z$ captures unlabeled nuisance variance. The
recognition model $q(y, z \mid x) = q(z \mid x, y) \, q(y \mid x)$ uses neural
network encoders, allowing nonlinear latent structure.

This framework is a natural generalization of DAS: DAS constrains the causal
variable to live on $\Gr(k,d)$ (a linear subspace), while the structured VAE
allows it to occupy a nonlinear manifold in activation space. The key advantage
is disentanglement: the supervised factor $y$ is trained to predict the
operation's output, while the unsupervised factor $z$ absorbs everything else.
If the causal variable is linear, the VAE encoder should learn a rotation
equivalent to DAS. If the causal variable is nonlinear (e.g., circular
structure from Fourier representations), the VAE can recover it where DAS
cannot.

The loss function combines the evidence lower bound (ELBO) with a supervised
classification term:
\begin{equation}
\mathcal{L} = \underbrace{-\mathbb{E}_{q}[\log p(x \mid y, z)]}_{\text{reconstruction}}
+ \beta \underbrace{D_\mathrm{KL}(q(y,z \mid x) \,\|\, p(y) p(z))}_{\text{regularization}}
+ \alpha \underbrace{\mathcal{L}_\text{CE}(y, \hat{y})}_{\text{supervision}}.
\label{eq:vae_loss}
\end{equation}

The encoder weights projecting to $y$ provide a linearized approximation of the
causal subspace; if this linearization has high overlap with the DAS subspace,
it confirms that the causal variable is approximately Grassmannian.

\subsection{Identifiability of the causal latent factor}
\label{sec:identifiability}

A central concern for any latent-variable method is \emph{identifiability}:
does the encoder recover the \emph{true} latent factors, or some
reparameterization that achieves good reconstruction without causal meaning?
We show that the Structured pi-SAE satisfies the conditions of iVAE
\citep{khemakhem2020variational}, which gives a formal guarantee that
$z_\text{causal}$ recovers the true causal variable up to a
component-wise transformation.

\begin{assumption}[Sufficient conditions for iVAE identifiability]
\label{ass:ivae}
\leavevmode
\begin{enumerate}[leftmargin=*,label=(\roman*)]
    \item \textbf{Auxiliary variable.} There exists an observed auxiliary
    variable $u$ (here, the task label $y$) such that the true latent prior
    factors as $p(z \mid u)$, with $p(z)$ not identifying $z$.
    \item \textbf{Label-conditional prior.} The conditional prior
    $p(z_\text{causal} \mid y)$ is an exponential family distribution whose
    sufficient statistics $T_j(z_j)$ are differentiable and whose natural
    parameters $\lambda_j(y)$ vary sufficiently across labels: for any
    $k$ latent dimensions, there exist $k{+}1$ labels $y_0, \ldots, y_k$
    such that the matrix $(\lambda(y_1) - \lambda(y_0), \ldots,
    \lambda(y_k) - \lambda(y_0))$ has rank $k$.
    \item \textbf{Injective mixing function.} The decoder
    $g_\theta: \mathcal{Z} \to \mathcal{H}$ is injective and
    differentiable.
\end{enumerate}
\end{assumption}

\begin{proposition}[Causal variable recovery]
\label{prop:recovery}
Suppose the Structured pi-SAE is trained to a global optimum of
$\mathcal{L}$ \emph{(Eq.~\ref{eq:vae_loss})} with $\beta > 0$ and
$\alpha > 0$, and Assumption~\ref{ass:ivae} holds.
Then the learned encoder $\hat{q}(z_\text{causal} \mid h, y)$ recovers the
true causal variable up to a component-wise reparameterization:
\begin{equation}
  \hat{z}_\text{causal} = \phi(z_\text{causal}),
  \label{eq:recovery}
\end{equation}
where $\phi: \R^k \to \R^k$ is a bijection acting
independently on each component ($\phi_j$ depends only on $z_j$).
\end{proposition}

\begin{proof}[Proof sketch]
At a global ELBO optimum, $q(z \mid h, y) = p(z \mid h, y)$ (the
variational gap vanishes), and the KL term forces the encoder to match
the label-conditional prior $p(z_\text{causal} \mid y)$.
Theorem~1 of \citet{khemakhem2020variational} then applies: any two encoders
achieving the same ELBO are related by a component-wise bijection on
$z_\text{causal}$, under Assumptions (i)--(iii).
The supervised term ($\alpha\mathcal{L}_\text{CE}$) additionally pins the
label-predictive direction, resolving the permutation ambiguity. \qed
\end{proof}

\paragraph{Checking the assumptions in practice.}
Assumption~(i) is satisfied by construction: the task label $y$ is
observed and the structured prior $p(z_\text{causal} \mid y)$ is
parameterized as $\mathcal{N}(\mu_y, \sigma_y^2 I)$ with per-label means
$\mu_y$.  This is an exponential family with sufficient statistics
$T_j(z_j) = (z_j, z_j^2)$ and natural parameters
$\lambda_j(y) = (\mu_{y,j}/\sigma_y^2,\, -1/(2\sigma_y^2))$, satisfying
Assumption~(ii) whenever per-label means are distinct (which holds
empirically: modular arithmetic outputs are uniformly distributed over
$\Z_p$, so the $k{+}1$ label pairs required for the rank condition
always exist).  Assumption~(iii) is enforced by the ELBO reconstruction
loss, which requires $g_\theta$ to approximately invert the encoder.

\paragraph{Three consequences.}
\textbf{(1) NL-DAS violates Assumption~(iii).}
Its decoder is trained exclusively on interchange loss with no reconstruction
constraint, so $g$ need not be injective.  Multiple latent codes may map to
the same activation, meaning the encoder's output is not uniquely determined
by the activation---and the iVAE guarantee does not apply
(\S\ref{sec:nldas_vacuous}).
\textbf{(2) The component-wise ambiguity is harmless for IIA.}
Interchange interventions swap $z_\text{causal}$ wholesale; $\phi$ is undone
by the decoder, leaving IIA invariant.
\textbf{(3) Linear DAS is the special case $\phi \in O(k)$.}
When the true causal variable is linear, the structured encoder degenerates
to a rotation, matching the DAS solution.

\subsection{Grokking as a phase transition}

Grokking is a training regime in which models memorize the training set long
before generalizing \citep{power2022grokking}. Modular arithmetic models trained
with weight decay exhibit grokking reliably for some operations but not others
\citep{nanda2023progress, zhong2024clock}. Prior work analyzed grokking through
Fourier representations and circular structure in the embedding space. Our
perspective is complementary: we ask how the causal subspace recovered by DAS
changes across the memorization-to-generalization transition, and whether that
transition is deterministic or stochastic across random seeds.
\citet{liu2023grokking} frame grokking as compression---moving from a memorized
high-complexity solution to a generalizing low-complexity one; our Grassmannian
perspective gives this compression a geometric interpretation as convergence to a
specific point on $\Gr(k,d)$.


\section{Methods}
\label{sec:methods}

\subsection{Operations and models}

We study 14 binary operations over $\Z_p$ (Table~\ref{tab:operations}), chosen
to span group actions, polynomials, piecewise functions, and non-algebraic maps.
For each operation, we train a single-layer transformer with
$d_\text{model} = 128$, four attention heads ($d_\text{head} = 32$), and
$d_\text{MLP} = 512$ on examples of the form $(a, b, {=}) \mapsto y$. Training
uses the standard grokking setup: 30\% train split, AdamW with weight decay 1.0,
and 15{,}000--80{,}000 epochs depending on the operation. We classify a model as
\emph{grokked} if its final test cross-entropy is below 0.1.

\begin{table}[t]
\centering
\small
\begin{tabular}{llll}
\toprule
Operation & Formula & Category & Seeds tested \\
\midrule
Multiplication & $ab \bmod 113$ & Group action & 42, 2024 \\
Comp.\ addition & $(a{+}b) \bmod 91$ & Group action & orig, 42 \\
Subtraction & $(a{-}b) \bmod 113$ & Group action & 42 \\
Division & $a/b \bmod 113$ & Group action & 42 \\
Bitwise XOR & $a \oplus b \bmod 113$ & Group action & 42, 137, 2024 \\
\midrule
Sum of squares & $(a^2{+}b^2) \bmod 113$ & Polynomial & 42 \\
Cubic sum & $(a^3{+}b^3) \bmod 113$ & Polynomial & 42 \\
Cubing & $a^3 \bmod 113$ & Polynomial & 42 \\
Squaring & $a^2 \bmod 113$ & Polynomial & 42 \\
Polynomial & $(a^2{+}b) \bmod 113$ & Polynomial & 42 \\
Affine & $(2a{+}3b{+}5) \bmod 113$ & Polynomial & 42 \\
\midrule
Max & $\max(a,b) \bmod 113$ & Piecewise & 42 \\
Abs.\ difference & $|a{-}b| \bmod 113$ & Piecewise & 42 \\
Power & $a^b \bmod 113$ & Non-algebraic & 42, 137 \\
\bottomrule
\end{tabular}
\caption{Operation atlas. Group actions provide clean equivariance targets.
Polynomials test whether nonlinear maps admit linear causal variables. Piecewise
and non-algebraic operations probe the boundary.}
\label{tab:operations}
\end{table}

\subsection{DAS fitting and $k$-sweeps}

For each trained model, we cache residual-stream activations and train a DAS
rotation $Q \in \R^{d \times k}$ using interchange interventions for 400
gradient steps. We sweep $k \in \{2, 4, 6, 8, 10, 12, 16, 20, 24, 32\}$ to
estimate the intrinsic dimension $k^*$, defined as the smallest $k$ achieving
$\IIA \geq 0.95$. Each $k$ value is an independent DAS optimization, not a
nested subspace---$k$-sweep profiles reveal whether the causal variable is
inherently low-dimensional or requires many dimensions.

\subsection{Hard-example selection for IIA}
\label{sec:hard_examples}

Standard IIA evaluation includes many ``easy'' examples where the intervention
does not change the model's top-1 prediction, inflating IIA regardless of whether
the subspace is causally valid. We define a \emph{hard example} as a
(base, source) pair where: (1) the base and source have different correct outputs
$y_b \neq y_s$, (2) the model's prediction matches the correct output for both
base and source separately, and (3) the interchange intervention at the DAS
subspace \emph{can} flip the model's prediction from $y_b$ to $y_s$. Hard IIA
measures the fraction of hard examples where the flip actually occurs. This
selection follows the methodology of CPCA-init DAS
(Tower et al., forthcoming), which showed that standard IIA saturates too easily
while hard IIA reveals genuine causal structure.

\subsection{Equivariance testing}

For operations with a natural group action, we test whether the DAS subspace
respects the algebraic symmetry. Given input $(a, b)$ with DAS projection
$z = Q^\top h(a,b)$, we construct a shifted input $(a{+}1 \bmod p, b)$ and
check whether the DAS projection rotates by $2\pi/p$ radians. The equivariant
fraction is the proportion of test inputs satisfying this rotation property.
Random orthogonal subspaces of the same dimension serve as controls; across all
operations, random-subspace equivariance is below 1\%.

\subsection{Structured VAE for nonlinear causal variables}
\label{sec:vae_method}

For each trained grokking model, we cache residual-stream activations at the
post-MLP hook and train a structured VAE with:
\begin{itemize}[leftmargin=*]
\item $z_\text{causal}$: $k$-dimensional latent factor ($k \in \{2,4,8\}$),
  semi-supervised with the operation's output label via a classification head.
\item $z_\text{nuisance}$: $m$-dimensional unsupervised factor ($m = 15$) to
  absorb non-causal variance.
\item Encoder: 2-layer MLP ($d_\text{model} \to 128 \to z$) with ReLU.
\item Decoder: 2-layer MLP ($z \to 128 \to d_\text{model}$) with ReLU.
\end{itemize}

Training uses the combined ELBO + classification loss (Eq.~\ref{eq:vae_loss})
for 300 epochs with $\beta = 1.0$ and $\alpha = 10.0$.

\paragraph{VAE-based IIA.}
We evaluate the VAE's causal subspace by performing interchange interventions in
the \emph{latent} space: given base activation $h_b$ and source activation
$h_s$, we encode both, swap $z_\text{causal}$ while keeping $z_\text{nuisance}$
from the base, decode back to activation space, and measure whether the model
produces the source's output. This is the nonlinear analogue of
Eq.~\ref{eq:intervention}: instead of $QQ^\top$ projection, we use the
encode-swap-decode path.

\paragraph{VAE equivariance.}
We test whether the VAE's $z_\text{causal}$ exhibits equivariance under additive
shifts, using the same protocol as for DAS (shift input $a \to a+1$, check
whether $z_\text{causal}$ rotates consistently). Because the VAE encoder is
nonlinear, it can represent circular structure directly in a 2D latent space.

\subsection{Nonlinear DAS baseline}
\label{sec:nldas}

To isolate the contribution of the VAE's generative structure from the benefit
of nonlinearity alone, we compare against an unconstrained nonlinear featurizer.
Given MLP encoder $f: \R^d \to \R^d$ and decoder $g: \R^d \to \R^d$, the
nonlinear DAS (NL-DAS) intervention is:
\begin{equation}
h' = g\bigl(f(h_b) - QQ^\top f(h_b) + QQ^\top f(h_s)\bigr)
\label{eq:nldas}
\end{equation}
where $Q \in \R^{d \times k}$ is learned jointly with $f$ and $g$ by maximizing
interchange accuracy. Both $f$ and $g$ use the same architecture as the VAE
encoder (2-layer MLP, 256 hidden units, ReLU), matching the VAE's representational
capacity. Crucially, NL-DAS has \emph{no reconstruction constraint}: $f$ and $g$
are trained exclusively on interchange loss and need not be (approximately)
inverse to each other.

We also evaluate NL-DAS+recon, which adds a reconstruction penalty
$\lambda \| g(f(h)) - h \|^2$ to the training loss, forcing $f$ and $g$ to
approximate mutual inverses. This tests whether the unconstrained NL-DAS result
degrades when the encoder-decoder cannot be degenerate.

\subsection{Intervention faithfulness metrics}
\label{sec:faithfulness_metrics}

Binary IIA measures \emph{interchange effectiveness}: does swapping the
subspace component change the model's output to match the source? But
effectiveness alone does not imply that the intervention is \emph{faithful}---that
it modifies only the targeted causal variable while preserving all other
information. A method that achieves IIA~=~1.0 by overwriting the entire
activation with a class-specific template is effective but not faithful.

We introduce four complementary metrics to evaluate intervention faithfulness:

\paragraph{Continuous distributional metrics.}
Binary IIA saturates at 1.0 whenever the top-1 prediction flips, hiding
differences in intervention quality. We measure the full distributional impact:
\begin{itemize}[leftmargin=*]
\item \textbf{KL divergence}: $\KL(p_\text{clean} \| p_\text{iv})$, where
  $p_\text{clean}$ and $p_\text{iv}$ are the softmax distributions before and
  after intervention. Lower KL means the intervention preserves more of the
  clean distribution beyond the swapped variable.
\item \textbf{Jensen-Shannon divergence}: the symmetric alternative
  $\JS = \frac{1}{2}\KL(p \| m) + \frac{1}{2}\KL(q \| m)$ where $m = (p+q)/2$.
  Bounded in $[0, \ln 2]$, more stable than KL when distributions have disjoint
  support.
\item \textbf{Probability difference}: $p_\text{iv}(y_s) - p_\text{iv}(y_b)$,
  a continuous analog of binary IIA that captures the model's confidence margin,
  not just top-1 agreement.
\item \textbf{Normalized logit difference}:
  $\frac{l_\text{iv}(y_s) - l_\text{iv}(y_b)}{|l_\text{clean}(y_s) - l_\text{clean}(y_b)|}$,
  where $l$ denotes logits. A value of 1.0 means the intervention reproduces the
  clean model's margin exactly; values $\gg 1$ or $\ll 0$ indicate distortion.
\end{itemize}

\paragraph{Intervention diversity.}
For each source label $y_s$, we compute the standard deviation of the intervened
activations $h'$ across all base inputs with the same source. If the decoder
produces the same $h'$ regardless of which base input was used---varying only
with $y_s$---the intervention is a lookup table: it does not preserve
base-specific nuisance information. We define the \emph{diversity ratio}:
\begin{equation}
\rho = \frac{\mathbb{E}_{y_s}\bigl[\text{std}(h'_{y_s})\bigr]}
            {\mathbb{E}_{y_s}\bigl[\text{std}(h_{y_s})\bigr]}
\label{eq:diversity_ratio}
\end{equation}
where the numerator is over intervened activations sharing source label $y_s$
and the denominator is over the corresponding natural source activations.
$\rho \approx 0$ indicates a lookup table (all intervened activations for the
same $y_s$ collapse to a single point). $\rho \approx 1$ indicates the
intervention preserves the natural variation within each class.

\paragraph{Reconstruction fidelity.}
For encoder-decoder methods (NL-DAS, VAE), we measure $\|g(f(h)) - h\|^2$ on
held-out activations. A high reconstruction MSE with high IIA is a warning sign:
the decoder may be ``hallucinating'' the correct output rather than faithfully
reconstructing the activation with a modified causal component.
The generative constraint is not merely a regularizer---it is a necessary
condition for identifiability (Proposition~\ref{prop:recovery}), as shown
empirically by pi-VAE's failure despite having the correct prior direction
(Table~\ref{tab:ablation}).

\subsection{Cross-distribution generalization}
\label{sec:cross_distribution}

Standard train/test splits draw from the same distribution. For IOI on GPT-2, we
additionally evaluate under distribution shift: (1)~\textbf{disjoint names}---train
on one set of proper names, evaluate on a completely disjoint set; and
(2)~\textbf{disjoint templates}---train on sentence templates 1--3, evaluate on
templates 4--5 plus novel hard templates. A method that generalizes under
distribution shift is finding a genuine causal variable; one that fails is
exploiting distributional regularities in the training set.


\section{Results}
\label{sec:results}

\subsection{The atlas: a three-class partition}

The 14 operations partition cleanly into three classes based on whether DAS
recovers a Grassmannian causal variable (Table~\ref{tab:main_results}).

\begin{table}[t]
\centering
\small
\setlength{\tabcolsep}{3pt}
\begin{tabular}{lcccccl}
\toprule
Operation & Grok & Loss & IIA($k{=}2$) & $k^*$ & Equiv. & Class \\
\midrule
\multicolumn{7}{l}{\textit{Always Grassmannian}} \\
Multiplication & \checkmark & 0.000 & 1.000 & 2 & 0.995 & Always \\
Subtraction & \checkmark & 0.000 & 1.000 & 2 & 1.000 & Always \\
Division & \checkmark & 0.000 & 1.000 & 2 & 1.000 & Always \\
Bitwise XOR & \checkmark & 0.003 & 1.000 & 2 & 1.000 & Always \\
Cubic sum & \checkmark & 0.000 & 1.000 & 2 & 1.000 & Always \\
Sum of squares & \checkmark & 0.005 & 1.000 & 2 & 0.987 & Always \\
Max & \checkmark & 0.074 & 0.960 & 2 & 0.957 & Always \\
\midrule
\multicolumn{7}{l}{\textit{Stochastic}} \\
Comp.\ addition$^*$ & \checkmark & 0.000 & --- & --- & 0.998 & Stochastic \\
Comp.\ addition (s42) & $\times$ & 10.207 & 0.800 & 6 & 0.110 & Stochastic \\
Power (s137) & \checkmark & 0.088 & 0.940 & 4 & 0.965 & Stochastic \\
Power (s42) & $\times$ & 1.143 & 0.980 & 2 & 0.665 & Stochastic \\
\midrule
\multicolumn{7}{l}{\textit{Never Grassmannian}} \\
Cubing & $\times$ & 9.765 & 0.857 & 4 & 0.509 & Never \\
Squaring & $\times$ & 7.808 & 0.857 & 4 & 0.314 & Never \\
Abs.\ difference & $\times$ & 1.835 & 0.800 & $>$32 & 0.189 & Never \\
Polynomial & $\times$ & 21.550 & 0.720 & $>$32 & 0.015 & Never \\
Affine & $\times$ & 26.209 & 0.780 & $>$32 & 0.026 & Never \\
\bottomrule
\end{tabular}
\caption{Atlas results for 14 modular arithmetic operations. $k^*$ is the
smallest $k$ achieving $\IIA \geq 0.95$. The ``Always'' class contains operations
that grok and produce equivariance $>95\%$ across all tested seeds.
$^*$Original seed (unknown value); k-sweep not recorded.}
\label{tab:main_results}
\end{table}

\subsection{Linear DAS fails on grokked modular addition}
\label{sec:linear_das_fails}

The starkest evidence that the Grassmannian assumption can fail even for correct,
generalizing models comes from modular addition. We train DAS at
$k \in \{2, 4, 8, 16\}$ on residual-stream activations from a fully grokked
model (test loss $< 10^{-7}$, test accuracy 100\%). At \emph{every} subspace
dimension, DAS returns $\IIA = 0.0$---not low, but exactly zero.

This is not a failure of DAS optimization. The model has learned the correct
modular addition algorithm via Fourier representations \citep{nanda2023progress}:
the ``identity'' of each number $a$ is encoded as
$(\cos 2\pi f a / p, \sin 2\pi f a / p)$ for key frequencies $f$. This
representation lives on $S^1$ (a circle embedded in $\R^d$), not on a linear
subspace of $\R^d$. DAS searches $\Gr(k,d)$ for a flat plane that captures a
curved manifold---a category error, not an optimization failure.

This motivates the Structured pi-SAE approach (\S\ref{sec:vae_results}):
a nonlinear sparse encoder can map the circular representation to a structured
latent space where interchange interventions succeed. On grokked addition, the
Structured pi-SAE achieves IIA~=~1.0 at $k{=}2$, confirming that the causal
variable exists but is nonlinear.

\subsection{Stochastic grokking: the same operation, opposite outcomes}
\label{sec:stochastic}

The stochastic class provides the strongest evidence for our central claim.
Power ($a^b \bmod 113$) grokked at seed 137 (test loss 0.088, equivariance
96.5\%) but not at seed 42 (test loss 1.14, equivariance 66.5\%). Composite
addition ($a + b \bmod 91$) grokked at one seed (test loss 0.000, equivariance
99.8\%) but not at seed 42 (test loss 10.2, equivariance 11.0\%). Same
operation, same architecture, same hyperparameters---opposite outcomes from
random initialization alone.

% TODO: 10-seed stability results here

\subsection{Memorization artifacts: high IIA, no structure}
\label{sec:memorization}

Squaring and cubing achieve $\IIA = 0.86$ at $k = 2$ and $\IIA = 1.0$ by $k = 4$
without grokking (test losses 7.8 and 9.8). Equivariance exposes the artifact:
31.4\% for squaring and 50.9\% for cubing, far below the $>$95\% threshold. The
lesson: \textbf{high IIA with low-$k$ saturation is compatible with a lookup
table}. ``Works for interchange'' and ``is a Grassmannian causal variable'' are
different claims.

\subsection{$k$-sweep profiles correlate with equivariance}

Grassmannian variables saturate at $\IIA \approx 1.0$ by $k = 4$: their causal
structure is inherently 2-dimensional. Non-Grassmannian variables show gradual IIA
growth that may not reach 1.0 even at $k = 32$. Polynomial reaches
$\IIA = 0.72$ at $k = 2$ and requires $k \geq 8$ to approach 0.90.

\subsection{Structured pi-SAE recovers nonlinear causal variables}
\label{sec:vae_results}

We upgrade the vanilla structured VAE to a \textbf{Structured pi-SAE}: a
label-conditional sparse autoencoder that partitions the latent space into
$z_\text{causal}$ (supervised, $L_1$-penalized) and $z_\text{nuisance}$
(unsupervised). The ``pi'' denotes the label-conditional prior from
\citet{narayanaswamy2017learning}; the ``SAE'' denotes $L_1$ sparsity on
$z_\text{causal}$ with an expansion factor of $8\times$ (so $k$ causal
dimensions expand to $8k$ sparse features). This combines the identifiability
guarantees of auxiliary supervision \citep{khemakhem2020variational} with the
interpretability benefits of sparse dictionary learning
\citep{bricken2023towards}.

\paragraph{The $2 \times 2$ ablation.}
Both components---structured prior and sparsity---are necessary
(Table~\ref{tab:ablation}). On grokking tasks: (1)~the plain SAE (no
label-conditional prior) achieves IIA~=~1.0 on addition but degrades on harder
operations (multiplication~0.72, quartic sum~0.32); (2)~the pi-VAE (structured
prior, no sparsity) achieves IIA~$\approx$~0 across all operations;
(3)~\emph{only} the Structured pi-SAE achieves IIA~=~1.0 uniformly across all
grokked operations and IIA~$\approx$~0 on all non-grokked operations.

\begin{table}[t]
\centering
\small
\setlength{\tabcolsep}{4pt}
\begin{tabular}{@{}l|cc|cc@{}}
\toprule
& \multicolumn{2}{c|}{\textbf{Plain prior} $\mathcal{N}(0,I)$}
& \multicolumn{2}{c}{\textbf{Structured prior} $p(z \mid y)$} \\
\textbf{Operation} & VAE & SAE & pi-VAE & \textbf{Str.\ pi-SAE} \\
\midrule
\multicolumn{5}{l}{\textit{Grokked}} \\
Addition       & 0.02 & 1.00 & 0.00 & \textbf{1.00} \\
Multiplication & 0.00 & 0.72 & 0.00 & \textbf{1.00} \\
Quartic sum    & 0.02 & 0.32 & 0.00 & \textbf{1.00} \\
IOI (GPT-2)    & 0.56 & 0.83 & 0.95 & \textbf{0.98} \\
\midrule
\multicolumn{5}{l}{\textit{Non-grokked (correct answer: $\approx 0$)}} \\
Mixed product  & 0.03 & 0.01 & 0.01 & 0.04 \\
Squaring       & 0.00 & 0.00 & 0.00 & 0.00 \\
\bottomrule
\end{tabular}
\caption{$2{\times}2$ ablation: IIA at $k{=}2$.
Plain SAE works on addition but degrades on harder operations (multiplication
$0.72$, quartic sum $0.32$) because without the structured prior, the sparse
encoder has no target direction.  pi-VAE (structured prior, no sparsity)
achieves $0.00$ on grokked operations---the prior provides the right
direction but without sparsity the encoder spreads the causal signal across
all dimensions, which cannot be cleanly swapped.  Only Structured pi-SAE
achieves $1.00$ uniformly on all grokked operations and $\leq 0.04$ on all
non-grokked operations.}
\label{tab:ablation}
\end{table}

\paragraph{Reconstruction MSE as a falsifiability criterion.}
For non-grokked operations, reconstruction MSE is $11$--$24$; for grokked
operations, $0.6$--$1.2$.  High MSE with near-zero IIA means the model has
\emph{no structured causal variable to recover}---not that the method
failed.  This diagnostic is unavailable for NL-DAS, which returns IIA $=
0.6$--$0.8$ on non-grokked operations (false positives) because it has no
reconstruction constraint.

\paragraph{End-to-end intervention training.}
For language model tasks where the causal structure is more complex, we add an
\textbf{end-to-end (E2E) intervention CE loss}: after swapping $z_\text{causal}$
and decoding, we run the intervened activation through the remaining layers and
compute $-\log p(y_s \mid h')$. This differentiable objective directly
optimizes for intervention success, analogous to how DAS trains on interchange
accuracy. Training proceeds in two phases: 200 epochs of reconstruction +
classification warmup, followed by E2E intervention training with an additive
intervention $h' = h_b + \text{dec}(z_\text{swap}) - \text{dec}(z_\text{orig})$
that cancels reconstruction error.

\paragraph{Results on GPT-2 language tasks.}
We evaluate on five tasks from the MIB benchmark \citep{prakash2024mib} at
layer~8 of GPT-2-small: gender bias, hypernymy, greater-than, SVA, and
capitals. Table~\ref{tab:nlp_results} reports the headline result.

\begin{table}[t]
\centering
\small
\setlength{\tabcolsep}{4pt}
\begin{tabular}{lcccc}
\toprule
\textbf{Method} & \multicolumn{2}{c}{\textbf{Hypernymy}} & \multicolumn{2}{c}{\textbf{Gender bias}} \\
& $k{=}1$ & $k{=}4$ & $k{=}1$ & $k{=}4$ \\
\midrule
DAS              & 0.07 & 0.23 & 0.62 & 0.88 \\
NL-DAS           & 0.72 & 0.92 & 1.00 & 1.00 \\
Str.\ pi-SAE     & 0.58 & 0.52 & 0.41 & 0.52 \\
\textbf{Str.\ pi-SAE E2E} & \textbf{0.97} & \textbf{0.98} & \textit{pending} & \textit{pending} \\
\bottomrule
\end{tabular}
\caption{Standard IIA on GPT-2 language tasks (layer~8, additive intervention).
Structured pi-SAE with E2E training achieves 0.97 standard IIA on hypernymy at
$k{=}1$, where DAS gets 0.07 and even NL-DAS gets only 0.72. Strict IIA
(argmax = source token) is 0.90 at $k{=}1$.}
\label{tab:nlp_results}
\end{table}

On hypernymy, the Structured pi-SAE with E2E training achieves standard
IIA~=~0.97 at $k{=}1$---$14\times$ better than DAS (0.07) and $1.35\times$
better than NL-DAS (0.72). Under the stricter argmax criterion (source token
must beat all 50{,}000 vocabulary tokens), IIA remains 0.90. The E2E objective
closes the gap between training and evaluation by directly optimizing the
quantity we measure.

On the Indirect Object Identification task \citep{wang2022interpretability},
DAS achieves IIA~=~0.30 at $k{=}2$, consistent with the IOI causal
variable spanning multiple attention heads.  Structured pi-SAE achieves
$0.98$.  NL-DAS achieves $1.00$, but the diversity ratio $\rho \approx 0$
(vs.\ $\rho \approx 0.91$ for the pi-SAE) exposes it as vacuous
(Table~\ref{tab:nldas_vacuous}).  The identifiability guarantee
(Proposition~\ref{prop:recovery}) applies: each IOI sentence has a distinct
indirect object label, the per-label means $\mu_y$ are well-separated, and
the ELBO reconstruction constraint forces approximate decoder injectivity.

\subsection{Unconstrained nonlinear DAS is vacuous}
\label{sec:nldas_vacuous}

A natural response to DAS's linearity constraint is to prepend an MLP
featurizer: learn a nonlinear coordinate system in which the causal variable
\emph{becomes} linear, then apply standard DAS. This approach---which we call
NL-DAS (\S\ref{sec:nldas})---achieves IIA~=~1.000 on IOI at $k = 1$, compared
to DAS's 0.193. This looks like a dramatic improvement. It is not.

\paragraph{The lookup-table failure mode.}
NL-DAS is trained end-to-end on interchange loss with no reconstruction
constraint. The encoder $f$ and decoder $g$ need not be approximate inverses,
and the training objective provides no incentive for them to be. This creates a
degenerate solution: $f$ maps every activation to a space where the first
coordinate encodes the class label and the remaining coordinates are arbitrary;
$g$ ignores the base-specific content entirely and produces a ``canonical
activation for class $y_s$'' regardless of which base input was used.

The intervention diversity ratio (\S\ref{sec:faithfulness_metrics}) exposes this
failure. Table~\ref{tab:nldas_vacuous} reports results for IOI and three
grokking operations across $k \in \{1, 2, 4\}$.

\begin{table}[t]
\centering
\small
\setlength{\tabcolsep}{3pt}
\begin{tabular}{lc rrrr rr}
\toprule
Task & $k$ & DAS & NL-DAS & NL-DAS+r & VAE & $\rho_\text{NL}$ & $\rho_\text{VAE}$ \\
\midrule
IOI & 1 & 0.193 & 1.000 & \textbf{---} & 0.990 & \textbf{---} & \textbf{---} \\
IOI & 2 & 0.297 & 1.000 & \textbf{---} & 0.980 & \textbf{---} & \textbf{---} \\
IOI & 4 & 0.556 & 1.000 & \textbf{---} & 0.978 & \textbf{---} & \textbf{---} \\
Addition & 1 & 0.006 & 0.194 & \textbf{---} & 0.000 & \textbf{---} & \textbf{---} \\
Addition & 4 & 0.289 & 0.883 & \textbf{---} & 0.089 & \textbf{---} & \textbf{---} \\
Mult.\ & 4 & 0.344 & 0.961 & \textbf{---} & 0.000 & \textbf{---} & \textbf{---} \\
Squaring & $*$ & 0.000 & 0.011 & \textbf{---} & 0.000 & \textbf{---} & \textbf{---} \\
\bottomrule
\end{tabular}
\caption{NL-DAS achieves high IIA by learning degenerate encoder-decoders.
  The VAE column reports Structured pi-SAE results.
  $\rho$ is the intervention diversity ratio (Eq.~\ref{eq:diversity_ratio});
  NL-DAS+r adds a reconstruction penalty.
  \textbf{---} marks values pending from hard-mode experiments.
  All controls (random labels, reconstruction-only, untrained, unconstrained
  plain VAE) achieve IIA~$< 0.09$, confirming the Structured pi-SAE is not
  cheating.}
\label{tab:nldas_vacuous}
\end{table}

Three lines of evidence confirm NL-DAS is vacuous:

\begin{enumerate}[leftmargin=*]
\item \textbf{Controls rule out the structured VAE cheating.} Random-label VAE
  (IIA $\leq 0.02$), reconstruction-only VAE with $\alpha = 0$
  (IIA $\leq 0.005$), untrained VAE (IIA = 0.000), and an unconstrained plain
  VAE with no causal/nuisance split (IIA $\leq 0.08$) all fail. The structured
  VAE's success requires correct supervision \emph{and} the
  causal-nuisance partition. This is the empirical counterpart of Proposition~\ref{prop:recovery}: NL-DAS
  fails Assumption~\ref{ass:ivae}(iii) because its decoder is not injective,
  so the iVAE identifiability guarantee does not apply
  \citep{khemakhem2020variational}. Without inductive biases, unsupervised
  disentanglement is impossible \citep{locatello2019challenging}; the
  structured prior and reconstruction constraint are not optional regularizers
  but necessary conditions.

\item \textbf{NL-DAS+recon degrades.} When forced to reconstruct its input
  ($g(f(h)) \approx h$), NL-DAS can no longer achieve degenerate solutions.
  % TODO: Fill with hard-mode results showing NL-DAS+recon IIA values

\item \textbf{Cross-distribution generalization.}
  % TODO: Fill with disjoint-name and disjoint-template results
  NL-DAS performance under distribution shift, where the encoder cannot
  memorize specific activation patterns from training, is pending from
  hard-mode experiments.
\end{enumerate}

\paragraph{Why the VAE is not vacuous.}
The structured VAE's ELBO provides three constraints that prevent the
lookup-table failure mode: (1)~the reconstruction loss forces the decoder to
produce activations close to the original, not class-specific templates;
(2)~the KL regularization prevents the encoder from collapsing the latent space
to a few discrete points; and (3)~the causal-nuisance partition ensures that
$z_\text{nuisance}$ absorbs base-specific information, so the decoder must use
it. Together, these constraints mean that swapping $z_\text{causal}$ while
preserving $z_\text{nuisance}$ genuinely modifies only the causal component.

\paragraph{Implications.}
Any nonlinear method for causal abstraction must be evaluated with intervention
faithfulness metrics, not just IIA. The NL-DAS failure mode is not specific to
our architecture---it applies to any unconstrained encoder-decoder trained
exclusively on interchange loss. Prior work proposing neural network featurizers
for causal abstraction \citep{geiger2024causal} should be re-evaluated with
diversity ratio and reconstruction checks.

\subsection{Continuous metrics reveal what binary IIA hides}
\label{sec:continuous_metrics}

Binary IIA treats a 51\% probability shift the same as a 99\% probability shift.
For the always-Grassmannian operations where IIA saturates at 1.0, continuous
metrics reveal important quality differences between methods.

% TODO: Insert continuous metrics table from hard-mode results
% showing KL, JS, prob_diff, normalized_logit_diff for DAS vs VAE vs NL-DAS

\subsection{Cross-distribution generalization on IOI}
\label{sec:cross_dist_results}

% TODO: Fill with disjoint-name and disjoint-template results from hard-mode
% Key question: does NL-DAS generalize to unseen names/templates, or does
% it exploit distributional regularities?

\subsection{Surprising positive and negative cases}

Cubic sum ($a^3 + b^3 \bmod 113$) achieves perfect equivariance (100\%) despite
cubing alone reaching only 50.9\%. The boundary is not ``group action versus
polynomial'' but whether the operation's binary structure provides sufficient
additive symmetry.

Affine ($2a + 3b + 5 \bmod 113$) is a linear function that \emph{fails} to
produce a Grassmannian variable. Under an additive shift $a \mapsto a + g$, the
output shifts by $2g$, not $g$. Equivariance requires the operation to be a group
action, not merely a linear function.

Max ($\max(a,b) \bmod 113$) achieves 95.7\% equivariance despite not being a
group action. It is piecewise equivariant: when the argmax is stable under the
shift, max behaves like a shifted identity.


\section{Discussion}
\label{sec:discussion}

\subsection{A hierarchy of causal validity}

Our results establish a four-level hierarchy for evaluating causal variable
claims, where each level adds a stronger requirement:

\begin{enumerate}
\item \textbf{Interchange effectiveness} (IIA; necessary, not sufficient): High
  IIA means the method supports interchange but does not distinguish genuine
  causal structure from memorization (squaring achieves $\IIA = 1.0$ at $k = 4$
  without grokking) or from degenerate encoder-decoders (NL-DAS achieves
  $\IIA = 1.0$ at $k = 1$ via lookup tables).

\item \textbf{Interchange faithfulness} (diversity ratio, reconstruction;
  necessary for nonlinear methods): A faithful intervention modifies the causal
  variable while preserving all other information. The diversity ratio
  $\rho$ (Eq.~\ref{eq:diversity_ratio}) and reconstruction MSE detect
  lookup-table failure modes. Without faithfulness, ``high IIA'' means only
  ``the decoder can produce the right output,'' not ``the encoder identified
  the right variable.'' Linear DAS is faithful by construction (the projection
  $QQ^\top$ preserves the orthogonal complement), but nonlinear methods must
  demonstrate faithfulness empirically.

\item \textbf{Distributional fidelity} (KL, JS, continuous metrics;
  distinguishes quality among effective methods): Among methods achieving
  IIA~$\approx$~1.0, continuous distributional metrics reveal how much of
  the clean distribution is preserved versus distorted. A method that flips
  the top-1 prediction (high IIA) while scrambling the rest of the distribution
  (high KL) is less valid than one that reproduces the full counterfactual
  distribution.

\item \textbf{Structural diagnostics} (equivariance, cross-distribution
  generalization): High equivariance means the recovered variable transforms
  coherently under the operation's symmetry group---the strongest evidence that
  the method has found a genuine causal variable rather than exploiting
  distributional artifacts. Cross-distribution generalization (disjoint names,
  novel templates) provides the analogous test for natural-language tasks that
  lack algebraic symmetries.
\end{enumerate}

The hierarchy has an important asymmetry. Linear DAS automatically satisfies
level~2 (faithfulness) because the orthogonal projection modifies only the
subspace component, but it may fail level~1 when the causal variable is
nonlinear. Unconstrained nonlinear methods easily satisfy level~1 but fail
level~2---they can achieve perfect IIA by learning degenerate solutions. The
structured VAE satisfies both because the ELBO provides the reconstruction and
regularization constraints that prevent degeneracy.

\subsection{Grokking enables causal structure---linear or nonlinear}

The stochastic grokking discovery shows that grokking is the process by which a
model's internal representation transitions from an unstructured memorization
encoding to a structured causal encoding. For operations with linear algebraic
structure, this structured encoding lives on $\Gr(k,d)$ and DAS recovers it.
For operations with nonlinear structure, the encoding lives on a curved manifold
and requires nonlinear methods to access.

The key insight is that grokking is about \emph{causal structure}, not
\emph{linear} causal structure. The Grassmannian perspective captures the linear
case precisely, but the broader phenomenon is the emergence of any structured
causal variable---linear or nonlinear---during the generalization transition.

\subsection{Implications for language model interpretability}

For practitioners applying DAS or nonlinear causal abstraction to language
models:
\begin{enumerate}
\item \textbf{Never report IIA alone.} IIA is necessary but insufficient at
  every level: linear DAS on memorized models (squaring, $\IIA = 1.0$) and
  unconstrained nonlinear methods (NL-DAS, $\IIA = 1.0$) both achieve
  perfect IIA vacuously. Always report at least one faithfulness metric
  (diversity ratio for nonlinear methods, equivariance or cross-distribution
  generalization for any method).
\item \textbf{Use continuous distributional metrics.} KL divergence and
  normalized logit difference distinguish among methods that all achieve
  $\IIA \approx 1.0$. Binary IIA treats a 51\% probability shift the same as
  99\%; continuous metrics capture this difference.
\item \textbf{Be suspicious of unconstrained nonlinear featurizers.} Any
  encoder-decoder trained end-to-end on interchange loss without reconstruction
  or regularization constraints can learn a lookup table. The ELBO
  (or an equivalent generative constraint) is necessary to prevent this
  failure mode.
\item \textbf{Test cross-distribution generalization.} For natural-language
  tasks, evaluate on disjoint entity sets or novel templates. A method that
  fails under distribution shift is exploiting surface regularities, not
  recovering causal structure.
\item \textbf{Consider nonlinear methods} when DAS IIA is low. Low IIA does not
  mean the causal variable does not exist---it may mean DAS is searching the
  wrong function class. But validate with faithfulness metrics.
\end{enumerate}

\subsection{Limitations}

Our experiments use single-layer transformers on synthetic tasks. Whether the
grokking--Grassmannian link holds in deeper models or natural language is an open
question, though the $k$-sweep and equivariance diagnostics apply regardless. The
structured VAE introduces a model-selection problem (architecture, $\beta$,
$\alpha$) that DAS avoids. Multi-seed stability experiments (10 seeds per
stochastic operation) are underway. The connection between the structured VAE
latent space and the Fourier representation structure identified by
\citet{nanda2023progress} deserves further theoretical analysis.


\section{Related Work}
\label{sec:related}

\paragraph{Causal abstraction.}
Causal abstraction provides the theoretical foundation for DAS
\citep{geiger2021causal, geiger2023finding}. \citet{wu2024interpretability}
scaled DAS via Boundless DAS. \citet{makelov2024illusion} identified
interpretability illusions where dormant pathways inflate IIA. Our geometric
diagnostics address a complementary problem: even within the linear intervention
class, high IIA does not imply genuine linear structure.

\paragraph{Grokking.}
Grokking was introduced by \citet{power2022grokking} and mechanistically
analyzed through Fourier representations \citep{nanda2023progress,
zhong2024clock}. Our contribution is to show that grokking has a specific
consequence for causal abstraction: it is the process that converts
non-Grassmannian representations into Grassmannian ones (for linear operations)
or into structured nonlinear ones (for nonlinear operations). Stochastic grokking
has been observed via the slingshot mechanism \citep{thilak2022slingshot} but not
connected to causal geometry.

\paragraph{Disentangled representations and causal variables.}
\citet{narayanaswamy2017learning} introduced semi-supervised VAEs for
disentanglement. Recent work on identifiable representations
\citep{khemakhem2020variational} establishes conditions under which latent
factors can be recovered. \citet{locatello2019challenging} proved that
unsupervised disentanglement is impossible without inductive biases; our
structured prior is precisely the auxiliary supervision their theorem
requires. Our application differs in that we evaluate
disentanglement through causal interchange (IIA), not reconstruction quality,
connecting the VAE literature to causal abstraction.

\paragraph{Representation geometry.}
Grassmannian methods have been applied to subspace tracking and domain adaptation
\citep{edelman1998geometry, gopalan2011domain}. Linear representation hypotheses
in interpretability \citep{marks2023geometry, engels2024language,
park2024linear} argue that many concepts are encoded in linear subspaces; our
work tests this hypothesis causally rather than correlationally.

\paragraph{Nonlinear causal abstraction.}
Concurrent work on nonlinear interchange interventions
\citep{geiger2024causal} uses neural network featurizers before applying DAS.
Our results show that this approach is vulnerable to the lookup-table failure
mode (\S\ref{sec:nldas_vacuous}): unconstrained featurizers achieve perfect IIA
by learning degenerate encoder-decoders. Our structured VAE approach uses the
ELBO to provide the reconstruction and regularization constraints that prevent
this degeneracy, connecting to the identifiability literature
\citep{khemakhem2020variational, sutter2020multimodal} which establishes that
auxiliary supervision is necessary for recovering latent factors. The NL-DAS
failure mode and our faithfulness metrics provide empirical evidence for this
theoretical requirement in the causal abstraction setting.


\section{Conclusion}
\label{sec:conclusion}

DAS searches for causal variables in a specific function class: linear subspaces
on the Grassmannian. Our atlas of 14 operations maps where this search succeeds
and where it fails. The boundary is sharp, governed by grokking, and stochastic
for operations near the transition.

Three practical implications follow. First, IIA should never be reported alone:
$\IIA = 1.000$ is compatible with a genuine causal variable, a memorized lookup
table, and a degenerate encoder-decoder. The four-level hierarchy we propose---interchange effectiveness, interchange faithfulness, distributional fidelity,
and structural diagnostics---provides a complete evaluation framework.

Second, nonlinear extensions of causal abstraction require generative structure
and structured sparsity. Unconstrained nonlinear featurizers achieve perfect IIA
vacuously; the Structured pi-SAE's combination of ELBO, label-conditional prior,
and $L_1$ sparsity prevents this failure mode. End-to-end intervention training
closes the gap between training and evaluation objectives, achieving IIA~=~0.97
on GPT-2 hypernymy where DAS gets 0.07. This connects causal abstraction to the
identifiability literature: auxiliary supervision plus generative constraints are
necessary for valid nonlinear causal variable recovery.

Third, continuous distributional metrics (KL, JS, normalized logit difference)
should replace binary IIA as the primary evaluation measure. Binary IIA
saturates too easily and hides important quality differences between methods.

The linearity assumption underlying DAS is not a benign default. It is a testable
hypothesis. And the natural fix---adding nonlinearity---introduces a new failure
mode that is invisible to the standard evaluation metric. This paper provides
the diagnostics to detect both problems and the structured nonlinear method
that avoids them.

\paragraph{Data and code.} All code for model training, DAS fitting, structured
VAE training, geometric diagnostics, and figure generation is available at
\url{[repository URL]}.

\bibliographystyle{tmlr}
\bibliography{references}

\end{document}
